\chapter{Thiết kế thực nghiệm cho phép biểu diễn SAT}
\chapterlead{Đánh giá trên bộ kiểm chuẩn không chỉ trả lời bộ biểu diễn nào nhanh
hơn. Một thực nghiệm tốt còn phải giải thích hiệu ứng đến từ kích thước, sức
lan truyền, bộ giải, đặc trưng hay cấu trúc của bài toán thử.}

\index{benchmark}\index{timeout}\index{PAR-2}
\index{performance profile}\index{cactus plot}
\index{virtual best solver}\index{shifted geometric mean}

\flowdiagram{Câu hỏi}{Ma trận thiết kế}{Bộ kiểm chuẩn}
  {Thước đo}{Phân tích}

\section{Thực nghiệm là một phép đo có mô hình}

Một bảng thời gian chạy không tự nói bộ biểu diễn nào tốt hơn. Quan sát được tạo
bởi cả
chuỗi
\[
  \text{bài toán thử}
  \rightarrow\text{mô hình}
  \rightarrow\text{bộ biểu diễn}
  \rightarrow\text{bộ tiền xử lý}
  \rightarrow\text{cấu hình bộ giải}
  \rightarrow\text{phần cứng}.
\]
Thay bất kỳ mắt xích nào cũng có thể đổi thứ hạng. Vì vậy kết luận cần nêu rõ
\emph{population} (quần thể mà kết luận hướng đến): một họ bài toán thử, một
miền tham số, một bộ giải hay một danh mục bộ giải. Cần phân biệt
\emph{internal validity} (hiệu lực nội tại: hiệu ứng thật sự do yếu tố
đang xét) với \emph{external validity} (hiệu lực khái quát: hiệu ứng
có còn giữ ngoài tập kiểm chuẩn hay không).

Trong nghiên cứu phép biểu diễn, ba loại bằng chứng bổ sung nhau:
\begin{itemize}
  \item \textbf{bằng chứng cơ chế}: công thức đếm, vết lan truyền và xung đột
        nhỏ giải thích vì sao hiệu ứng xuất hiện;
  \item \textbf{bằng chứng có kiểm soát}: \emph{ablation}
        (thí nghiệm phân tích thành phần (ablation study)) giữ cố định các yếu tố khác;
  \item \textbf{bằng chứng so sánh}: đánh giá trên nhiều họ bài toán và nhiều
        bộ giải để kiểm tra độ bền của kết luận.
\end{itemize}
Một kết quả có giá trị tham khảo cao cần ít nhất hai loại; thời gian chạy đơn
lẻ chỉ là bằng chứng so sánh yếu nếu thiếu cơ chế và biến kiểm soát.

\section{Bắt đầu từ một câu hỏi nhân quả}

Một nghiên cứu về phép biểu diễn thường có một trong bốn câu hỏi:
\begin{enumerate}
  \item phép biểu diễn mới giảm kích thước bao nhiêu;
  \item sức lan truyền thay đổi như thế nào;
  \item thời gian chạy và số bài giải được thay đổi ra sao;
  \item một mô-đun có tương tác với phá đối xứng, giải tăng dần hoặc bộ giải không.
\end{enumerate}

Câu hỏi càng rõ, ma trận thiết kế càng nhỏ. Nếu mục tiêu là đo bộ đếm dùng
chung, hãy giữ cố định bộ giải, tiền xử lý, chiến lược tìm kiếm và quy trình
xử lý hàm mục tiêu; chỉ thay mô-đun bộ đếm. Nếu đồng thời thay bộ giải và cách
phá đối xứng, kết quả không phân định được đóng góp riêng lẻ của từng thành phần.

\begin{designrule}{Đơn vị của thí nghiệm phân tích thành phần (ablation study)}
Một đặc trưng không có hiệu ứng độc lập khi làm thay đổi giao diện của phép biểu diễn.
Trong trường hợp đó, dùng cấu hình
\(\text{phép biểu diễn}\times\text{đặc trưng}\) làm đơn vị so sánh.
\end{designrule}

\section{Ma trận thiết kế}

Một ma trận điển hình gồm:
\[
  \text{bộ biểu diễn}
  \times
  \text{bộ giải}
  \times
  \text{đặc trưng}
  \times
  \text{bài toán thử}
  \times
  \text{random seed}.
\]
Ở đây \emph{random seed} (hạt giống ngẫu nhiên) là giá trị khởi tạo bộ sinh
số giả ngẫu nhiên; giữ lại giá trị này giúp tái lập đúng chuỗi lựa chọn ngẫu
nhiên của một lần chạy.
Không nhất thiết chạy tích Descartes đầy đủ. Có thể dùng hai giai đoạn:
\begin{enumerate}
  \item \textbf{sàng lọc}: chọn vài bộ giải đại diện và miền tham số rộng;
  \item \textbf{xác nhận}: giữ các cấu hình có triển vọng, tăng số bài toán
        thử và số hạt giống ngẫu nhiên.
\end{enumerate}

Mọi cấu hình phải dùng cùng ngữ nghĩa và bộ kiểm tra. Nếu một bộ biểu diễn có tiền
xử lý riêng, thời gian tiền xử lý phải được tính vào tổng thời gian chạy hoặc
được báo thành một cột riêng.

\section{Bộ kiểm chuẩn và phân tầng}

Một bộ kiểm chuẩn nên có siêu dữ liệu cấu trúc: số biến miền, mật độ ràng buộc,
độ rộng cửa sổ, ngưỡng, lớp đối xứng, khoảng \(LB\)--\(UB\), hoặc các đặc trưng
riêng của bài toán. Phân tầng theo những đặc trưng này giúp giải thích tại sao
một phép biểu diễn thắng trong một vùng.

Ba nguồn bài toán thử thường dùng:
\begin{itemize}
  \item bài toán chuẩn từ văn liệu;
  \item bài toán thực tế có nguồn gốc rõ ràng;
  \item bài toán sinh tổng hợp để quét có kiểm soát một tham số.
\end{itemize}

Các bài toán sinh tổng hợp cho phép quan sát điểm chuyển theo \(n,k,w\), mật độ
hoặc độ rộng miền; dữ liệu thực cung cấp bối cảnh sử dụng, còn bài toán chuẩn
tạo khả năng so sánh với kết quả trước. Báo riêng ba nguồn giúp người đọc hiểu
đúng vai trò của từng tập.

\section{Tách tinh chỉnh và đánh giá}

Tùy chọn bộ giải, độ rộng khối, kích thước nhóm và ngưỡng biểu diễn là các siêu
tham số. Nếu chọn chúng trên chính tập kiểm thử, thời gian chạy báo cáo sẽ lạc
quan quá mức. Một quy trình đơn giản:
\begin{enumerate}
  \item chia bộ kiểm chuẩn thành tập tinh chỉnh và tập đánh giá;
  \item khóa cấu hình sau khi tinh chỉnh;
  \item chỉ dùng tập đánh giá cho kết luận cuối;
  \item nếu dữ liệu ít, dùng \emph{cross-validation} (kiểm định chéo) theo họ
        thay vì chia ngẫu nhiên từng bài toán.
\end{enumerate}

Tinh chỉnh theo họ tránh việc các biến thể gần như trùng nhau xuất hiện ở cả
hai tập.

\section{Các thước đo}

\subsection{Kích thước và xây dựng}

Ghi ít nhất:
\[
  n_v,\quad n_c,\quad n_\ell,\quad
  \#\text{binary},\quad
  \#\text{ternary},\quad
  t_{\text{encode}},\quad
  \text{bộ nhớ cực đại}.
\]
Thời gian của bộ giải nên đo từ sau khi dựng công thức; đồng thời báo thời gian
toàn trình nếu bộ sinh công thức chiếm một phần đáng kể.

\subsection{Quá trình tìm kiếm}

Số quyết định, lượt lan truyền, xung đột, lần khởi động lại và mệnh đề đã học
giúp giải thích cơ chế tạo ra thời gian chạy quan sát được:
\[
  \text{lượt lan truyền/xung đột},\qquad
  \text{xung đột/giây},\qquad
  \text{bộ nhớ/biến}.
\]

\subsection{Right-censored runtime: thời gian chạy có kiểm duyệt}

Khi một lần chạy hết thời hạn, ta chỉ biết thời gian cần thiết lớn hơn giới hạn;
dữ liệu như vậy được gọi là \emph{right-censored} (dữ liệu bị giới hạn thời gian).
Vì vậy, thống kê thời gian chạy phải tính cả bài đã giải và bài hết thời hạn
bằng một trong các cách:
\begin{itemize}
  \item số bài giải được tại một giới hạn thời gian cố định;
  \item PAR-2 hoặc PAR-10, tức thời gian trung bình có phạt khi hết thời hạn;
  \item \emph{cactus plot} (biểu đồ Cactus (Cactus plot));
  \item \emph{performance profile} (hồ sơ hiệu năng);
  \item thời gian đạt mục tiêu cho bài toán tối ưu.
\end{itemize}

Hồ sơ hiệu năng của Dolan--Moré chuẩn hóa thời gian chạy theo cấu hình tốt nhất
trên từng bài toán thử. Với tỷ số
\[
  r_{p,s}=
  \frac{t_{p,s}}{\min_{s'}t_{p,s'}},
\]
đường
\[
  \rho_s(\tau)=
  \frac{1}{|P|}
  \card{\set{p\in P\mid r_{p,s}\le\tau}}
\]
cho biết tỷ lệ bài mà cấu hình \(s\) nằm trong hệ số \(\tau\) so với tốt nhất
\parencite{dolan2002profiles}. Giá trị tại \(\tau=1\) đo số lần thắng; phần
đuôi đo độ bền. Bài hết thời hạn phải được gán một tỷ số lớn nhất nhất quán và
được mô tả rõ.

\begin{keyidea}{Đọc số bài giải được và thời gian chạy cùng nhau}
Một cấu hình giải thêm các bài khó có thể có tổng thời gian chạy lớn hơn cấu
hình hết thời hạn sớm. Vì vậy không xếp hạng bằng một cột thời gian đơn lẻ.
\end{keyidea}

\begin{workedexample}{Hết thời hạn làm thay đổi cách xếp hạng}
Với giới hạn \(30\) giây, ta có:
\[
\begin{array}{c|ccc}
 & p_1&p_2&p_3\\\hline
\text{Đối chứng}&2&5&\text{hết hạn}\\
\text{Mới}&1&10&20
\end{array}
\]
Đối chứng nhanh hơn trên \(p_2\), nhưng chỉ giải \(2/3\) bài; cấu hình mới giải
cả ba. Nếu PAR-2 gán \(60\) giây cho lần chạy hết thời hạn:
\[
  \operatorname{PAR2}(\text{Đối chứng})=\frac{2+5+60}{3}=22.33,
  \qquad
  \operatorname{PAR2}(\text{Mới})=\frac{1+10+20}{3}=10.33.
\]
Trung bình chỉ trên bài đã giải sẽ cho đối chứng \(3.5\) giây và cấu hình mới
\(10.33\) giây, từ đó tạo kết luận ngược vì đã bỏ chính bài khó nhất.
\end{workedexample}

\begin{figure}[tb]
\centering
\begin{tikzpicture}[satfig,x=18mm,y=11mm]
  \draw[->,thick,Ink] (0,0)--(5.4,0) node[right,satlabel] {\(\tau\)};
  \draw[->,thick,Ink] (0,0)--(0,3.5) node[above,satlabel] {\(\rho_s(\tau)\)};
  \foreach \y/\lab in {1/0.33,2/0.67,3/1.00}{
    \draw[Rule] (0,\y)--(5,\y);
    \node[left=4pt,satlabel] at (0,\y) {\lab};
  }
  \draw[very thick,Blue]
    (0.5,0)--(0.5,1)--(1.2,1)--(1.2,2)--(3.8,2)--(3.8,3)--(5,3);
  \draw[very thick,Teal]
    (0.5,0)--(0.5,1)--(.9,1)--(.9,2)--(1.8,2)--(1.8,3)--(5,3);
  \node[satlabel,text=Blue] at (4.5,2.6) {Đối chứng};
  \node[satlabel,text=Teal] at (2.5,3.25) {Mới};
  \draw[Ink] (.5,2pt)--(.5,-2pt) node[below=3pt,satlabel]{1};
  \draw[Ink] (1.8,2pt)--(1.8,-2pt) node[below=3pt,satlabel]{2};
  \draw[Ink] (3.8,2pt)--(3.8,-2pt) node[below=3pt,satlabel]{4};
\end{tikzpicture}
\caption{Hồ sơ hiệu năng (performance profile): phần đầu cho thấy số bài giải nhanh nhất, phần đuôi
cho độ bền so với cấu hình tốt nhất trên từng bài toán thử.}
\label{fig:performance-profile}
\end{figure}

\subsection{Biểu đồ Cactus (Cactus plot) và biểu đồ phân tán}

Với một cấu hình \(s\), sắp tăng thời gian của các bài đã giải:
\[
  t_{(1),s}\le t_{(2),s}\le\cdots\le t_{(q_s),s}.
\]
Biểu đồ Cactus (Cactus plot) đặt số bài đã giải \(j\) trên trục ngang và
\(t_{(j),s}\) trên trục dọc. Đường đi xa hơn sang phải biểu thị độ bền; đường
thấp hơn biểu thị giải nhanh hơn ở cùng số bài đã giải. Đường dừng ở \(q_s\),
nên các bài hết thời hạn không bị âm thầm loại khỏi kết luận.

Biểu đồ phân tán dùng so sánh ghép cặp: mỗi điểm ứng với cùng một bài toán thử,
có tọa độ \((t_{p,A},t_{p,B})\), thường trên thang lôgarit. Điểm phía trên đường chéo
\(t_A=t_B\) nghĩa là \(A\) nhanh hơn; điểm phía dưới nghĩa là \(B\) nhanh hơn.
Bài hết thời hạn cần được đặt ở biên \(T\) bằng ký hiệu riêng hoặc mũi tên, và
SAT/UNSAT nên dùng dấu điểm khác nhau. Biểu đồ phân tán cho biết \emph{bài toán
nào} tạo khác biệt; biểu đồ Cactus (Cactus plot) cho biết \emph{toàn tập} được bao phủ
đến đâu.

\begin{figure}[tb]
\centering
\begin{tikzpicture}[satfig,x=8mm,y=8mm]
  \begin{scope}
    \draw[->,Ink] (0,0)--(5.2,0) node[right,satlabel]{số bài giải được};
    \draw[->,Ink] (0,0)--(0,4.2) node[above,satlabel]{thời gian};
    \draw[Blue,very thick] (.5,.35)--(1.3,.55)--(2.1,.9)--(2.9,1.45)--(3.7,2.4);
    \draw[Teal,very thick] (.5,.45)--(1.3,.7)--(2.1,1.2)--(2.9,2.1)--(4.5,3.5);
    \node[satlabel,text=Blue] at (3.55,2.05) {A};
    \node[satlabel,text=Teal] at (4.55,3.15) {B};
    \node[satlabel,font=\sffamily\bfseries\small] at (2.5,4.6) {Cactus plot};
  \end{scope}
  \begin{scope}[xshift=62mm]
    \draw[->,Ink] (0,0)--(4.6,0) node[right,satlabel]{\(t_A\)};
    \draw[->,Ink] (0,0)--(0,4.2) node[above,satlabel]{\(t_B\)};
    \draw[Rule,dashed] (.25,.25)--(4,4);
    \foreach \x/\y in {.6/1.2,1.1/.8,1.5/2.7,2.1/1.5,2.6/3.4,3.4/2.8}
      \fill[Blue] (\x,\y) circle (2.2pt);
    \draw[Gold,very thick] (4,2.4)--(4.35,2.4);
    \draw[Gold,very thick] (2.1,4)--(2.1,4.35);
    \node[satlabel,text=Blue,rotate=45] at (3.25,3.55) {\(t_A=t_B\)};
    \node[satlabel,font=\sffamily\bfseries\small] at (2.2,4.6) {Phân tán ghép cặp};
  \end{scope}
\end{tikzpicture}
\caption{Biểu đồ Cactus (Cactus plot) (trái) giữ số bài giải được trong hình; biểu đồ
phân tán (phải) giữ ghép cặp theo bài toán thử. Các dấu ở biên minh họa lần
chạy hết thời hạn của một cấu hình.}
\label{fig:cactus-scatter}
\end{figure}

\section{SAT và UNSAT là hai chế độ khác nhau}

Các bài toán SAT cần tìm một nghiệm làm chứng; các bài toán UNSAT cần đóng toàn bộ
không gian tìm kiếm. Một phép biểu diễn có thể mạnh ở SAT nhưng yếu ở UNSAT hoặc
ngược lại.
Bảng kết quả nên tách:
\[
  \#SAT,\quad \#UNSAT,\quad t_{SAT},\quad t_{UNSAT}.
\]

Trong tối ưu, các lần gọi bộ giải gần giá trị tối ưu thường có cả hai trạng
thái: mô hình tại cận khả thi và \UNSAT tại cận tốt hơn. Vì thế, hồ sơ theo
toàn bộ dãy cận chứa nhiều thông tin hơn thời gian chạy cuối cùng.

\section{Random seed và biến thiên giữa các lần chạy}

\emph{Random seed}, được gọi là hạt giống ngẫu nhiên sau lần giới thiệu đầu,
không phải dữ liệu của bài toán; tham số này là giá trị khởi tạo (seed) cho bộ sinh số giả ngẫu nhiên
của thuật toán. CDCL có yếu tố ngẫu nhiên trong phá hòa và tìm kiếm. Với bài
toán khó, cần chạy nhiều hạt giống ngẫu nhiên và báo trung vị, các phân vị hoặc
khoảng tin cậy. So sánh ghép cặp phải dùng cùng bài toán, cùng chỉ số hạt giống
ngẫu nhiên và cùng giới hạn thời gian, bộ nhớ.

Nếu bộ biểu diễn là tất định nhưng bộ giải có yếu tố ngẫu nhiên, hạt giống ngẫu
nhiên thuộc về bộ giải. Nếu bộ sinh dùng thứ tự ngẫu nhiên, cần tách hạt giống
ngẫu nhiên của bộ sinh khỏi hạt giống ngẫu nhiên của bộ giải.

\section{So sánh ghép cặp và độ lớn hiệu ứng}

Mỗi bài toán thử tạo một cặp quan sát tự nhiên giữa hai bộ biểu diễn. Vì phân bố
thời gian chạy thường lệch và có dữ liệu hết thời hạn, chỉ báo trung bình và
một kiểm định trên các lần chạy độc lập là chưa phù hợp. Nên báo đồng thời:
\begin{itemize}
  \item số lần thắng/hòa/thua trên cùng bài toán và cùng hạt giống ngẫu nhiên;
  \item trung vị của lôgarit hệ số tăng tốc
        \(\log(t_{\mathrm{base}}/t_{\mathrm{new}})\);
  \item khoảng bất định bằng \emph{bootstrap} (lấy mẫu lại) theo \emph{họ}, không chỉ
        theo từng bài toán;
  \item chênh lệch số bài giải được tại nhiều mốc thời gian.
\end{itemize}

Độ lớn hiệu ứng quan trọng hơn một giá trị \(p\) đứng riêng. Một bộ biểu diễn nhanh
hơn \(3\%\) trên hàng nghìn bài dễ nhưng làm mất năm bài khó có ý nghĩa khác
với bộ biểu diễn chỉ giảm ít trung vị nhưng mở rộng vùng giải được. Với tối ưu,
nên so cả thời gian đạt nghiệm tốt nhất đã biết, thời gian hoàn tất chứng minh
và khoảng cách tối ưu theo thời gian.

\section{Phân tích theo vùng cấu trúc}

Một kết luận toàn cục dễ che mất điểm chuyển ưu thế. Bộ đếm tuần tự có thể tốt
khi \(k/n\) nhỏ; mạng đếm có thể cải thiện khi cần nhiều ngưỡng; biểu diễn thứ tự
có thể thắng khi các khoảng cấm rộng. Vì vậy, trước khi vẽ kết quả, cần xác
định trước các tầng theo đặc trưng có ý nghĩa thuật toán:
\[
  \frac{k}{n},\quad
  \text{độ chồng lấn cửa sổ},\quad
  \text{độ rộng miền},\quad
  \text{mật độ ràng buộc},\quad
  \text{quy mô đối xứng}.
\]
Sau đó dùng bản đồ nhiệt hoặc bảng theo tầng để tìm \emph{miền ưu thế} thay vì
chỉ tìm ``người thắng''. Cách phân tích này cũng cung cấp dữ liệu cho bộ chọn
phép biểu diễn ở Chương 3.

\section{Kiểm đúng trước khi đo nhanh}

Trước khi chạy bộ kiểm chuẩn lớn:
\begin{enumerate}
  \item kiểm thử vét cạn trên các bài toán nhỏ;
  \item kiểm thử vi sai giữa hai bộ biểu diễn;
  \item giải mã và kiểm tra mọi mô hình SAT dùng trong bảng;
  \item tính lại hàm mục tiêu từ nghiệm, không đọc trực tiếp chi phí nội bộ;
  \item giữ một số công thức \CNF và chứng minh nhỏ làm phép thử hồi quy.
\end{enumerate}

Với \(n\) nhỏ, bộ liệt kê sinh mọi phép gán cho biến chính. So sánh phép chiếu
của các mô hình \CNF với tập nghiệm miền. Cách này phát hiện lỗi biên hiệu quả
hơn việc chỉ chạy các bài toán kiểm chuẩn lớn.

\section{Gói tái lập và khả năng truy vết}

\emph{Artifact} (gói phần mềm--dữ liệu tái lập) tối thiểu gồm:
\begin{itemize}
  \item mã nguồn và mã định danh bản sửa đổi chính xác;
  \item bản kê bài toán thử và mã kiểm tra toàn vẹn;
  \item môi trường cùng phiên bản bộ giải;
  \item lệnh chạy và giới hạn tài nguyên;
  \item nhật ký thô, bộ phân tích nhật ký và chương trình tạo bảng;
  \item bộ giải mã và bộ kiểm tra.
\end{itemize}
Gói tái lập tạo một đường đi kiểm tra được từ phát biểu đến bằng chứng thô.

\begin{figure}[tb]
\centering
\begin{tikzpicture}[satfig,node distance=8mm]
  \node[satbox] (claim) {Phát biểu\\``nhanh hơn''};
  \node[satbox,right=of claim] (table) {Bảng/biểu đồ};
  \node[satbox,right=of table] (parsed) {Nhật ký\\đã phân tích};
  \node[satbox,right=of parsed] (raw) {Nhật ký thô};
  \node[satbox,right=of raw] (run) {Lệnh chạy\\+ bản sửa đổi};
  \draw[satedge] (claim)--(table);
  \draw[satedge] (table)--(parsed);
  \draw[satedge] (parsed)--(raw);
  \draw[satedge] (raw)--(run);
  \node[satlabel,below=6mm of parsed]
    {Mỗi mũi tên phải có chương trình, mã kiểm tra hoặc siêu dữ liệu để truy vết.};
\end{tikzpicture}
\caption{Chuỗi truy vết tối thiểu của một phát biểu thực nghiệm.}
\label{fig:claim-artifact-chain}
\end{figure}

Số liệu thực nghiệm có thể được kiểm tra ở ba tầng:
\begin{enumerate}
  \item \textbf{báo cáo}: bài viết mô tả tập bài toán thử, giới hạn thời gian
        và các cấu hình đối chứng;
  \item \textbf{tái tạo}: mã nguồn, nhật ký và dữ liệu cho phép dựng lại các
        bảng và biểu đồ;
  \item \textbf{tái lập độc lập}: một môi trường khác thực hiện lại thí
        nghiệm theo cùng thiết kế.
\end{enumerate}
Tầng báo cáo hỗ trợ việc diễn giải kết quả; tầng tái tạo kiểm tra đường đi từ
nhật ký đến bảng; tầng tái lập độc lập đánh giá độ bền khi môi trường thực thi
thay đổi.

\section{Mẫu quy trình thực nghiệm}

\begin{summarybox}
\textbf{Câu hỏi.} Mô-đun nào được đánh giá?\\
\textbf{Giả thuyết.} Kỳ vọng tác động lên kích thước, sức lan truyền hay thời gian chạy?\\
\textbf{Yếu tố.} Bộ biểu diễn, bộ giải, đặc trưng, họ bài toán, hạt giống ngẫu nhiên.\\
\textbf{Biến kiểm soát.} Ngữ nghĩa, bộ kiểm tra, cận, tiền xử lý, phần cứng.\\
\textbf{Thước đo.} Kích thước, thời gian dựng, bộ đếm tìm kiếm, số bài giải được/PAR/hồ sơ.\\
\textbf{Kiểm đúng.} Kiểm thử vét cạn, kiểm thử vi sai và bộ kiểm tra mô hình.\\
\textbf{Đầu ra.} Nhật ký thô, biểu đồ, bảng tổng hợp và bản kê gói tái lập.
\end{summarybox}

\section{Bài học thiết kế}

\begin{enumerate}
  \item Đánh giá thực nghiệm bắt đầu từ một câu hỏi, không bắt đầu từ một bảng cấu hình.
  \item Tách tinh chỉnh khỏi đánh giá và phân tầng theo họ bài toán.
  \item Báo cả kích thước, bộ đếm tìm kiếm và thời gian chạy có dữ liệu hết hạn.
  \item Tách SAT/UNSAT và dùng nhiều hạt giống ngẫu nhiên ở vùng khó.
  \item Bộ kiểm tra mô hình là điều kiện trước của mọi bảng hiệu năng.
\end{enumerate}

\subsection*{Câu hỏi nghiên cứu}
\begin{enumerate}
  \item Thiết kế thí nghiệm phân tích thành phần (ablation study) \(2\times2\) cho bộ biểu diễn và phá đối xứng.
  \item Với ba lần hết thời hạn trong mười bài toán thử, tính PAR-2 và so với
        trung bình trên các bài đã giải.
  \item Lập một bản kê gói tái lập tối thiểu cho một thí nghiệm SAT tăng dần.
\end{enumerate}
